\documentclass[aps,prl,onecolumn,superscriptaddress,notitlepage]{revtex4-2}
\usepackage{graphicx}
\usepackage{amsmath,amssymb}
\usepackage[colorlinks=true,allcolors=blue]{hyperref}
\setlength{\tabcolsep}{7pt}
\renewcommand{\thefigure}{S\arabic{figure}}
\renewcommand{\thetable}{S\arabic{table}}
\renewcommand{\theequation}{S\arabic{equation}}
\begin{document}
\title{Supplemental Material for ``Sub-cycle rectification at the metal contact in antenna-enhanced
high-harmonic generation from solids''}
\author{[Author list to be completed]}
\maketitle

\section{Model}
All quantities are in atomic units unless stated otherwise. The coordinate $z$ runs along the field
line that leaves the metal; the metal occupies $-L_m<z<0$ and the crystal $d<z<d+L_c$, where $d$ is
the thickness of an optional insulating spacer. The static single-particle potential is
\begin{equation}
V(z)=\begin{cases}
V_m, & -L_m<z<0 \quad\text{(metal)},\\
E_c+h_s, & 0<z<d \quad\text{(spacer)},\\
-V_0\{1+\cos[2\pi(z-d)/a]\}, & d<z<d+L_c \quad\text{(crystal)},\\
E_c+\chi, & \text{elsewhere (vacuum)},
\end{cases}
\end{equation}
with all steps smoothed by $\tfrac12[1+\tanh(z/w)]$, $w=1.5$. The parameters are $a=8$ (0.423~nm),
$V_0=0.32$, which gives band edges $E_v=-1.05$~eV and $E_c=2.28$~eV relative to the maxima of the
crystal potential and a gap $E_g=3.34$~eV between the second and third bands; electron affinity
$\chi=4$~eV; spacer barrier $h_s=3.1$~eV; and $V_m=E_{\rm F}-5.5$~eV, where the Fermi level
$E_{\rm F}=E_c-\Phi_{\rm B}$ is set by the Schottky barrier. The reference geometry has
$L_m=240$ (12.7~nm) and $L_c=80a$ (33.9~nm). Each spatial orbital below $E_{\rm F}$ is doubly
occupied. The metal electron density is $2k_{\rm F}/(\pi a^2)=4.3\times10^{22}$~cm$^{-3}$
and its plasma frequency is 7.7~eV.

The time-dependent part of the Hamiltonian is $D(t)g(z)+U_{\rm H}(z,t)$. The first term is the
potential of an external charge sheet at the metal surface, which represents the oscillating
surface charge of the antenna, with
\begin{equation}
g(z)=\begin{cases}0,& z<0,\\ z/\varepsilon_s, & 0<z<d,\\ d/\varepsilon_s+\ell\,[1-e^{-(z-d)/\ell}], & z>d,\end{cases}
\end{equation}
where $\ell$ models the decay of the near field into the crystal and $\varepsilon_s=3.9$ is the
permittivity of the (passive) spacer. The driving waveform is
$D(t)=D_0\sin^2(\pi t/\tau)\cos[\omega(t-\tau/2)]$ for $0<t<\tau$. The induced Hartree potential
obeys the planar Poisson equation
\begin{equation}
\partial_z^2U_{\rm H}(z,t)=-4\pi\,\frac{2}{a^2}\Big[\sum_j|\psi_j(z,t)|^2-\sum_j|\psi_j(z,0)|^2\Big],
\end{equation}
with vanishing field far to the left; because the induced charge integrates to zero the field also
vanishes far to the right. The static potential plays the role of a frozen ground-state Kohn--Sham
potential, so the treatment corresponds to the time-dependent Hartree (random-phase) approximation
for the induced field; exchange-correlation and image-potential effects are not included.

\section{Numerical method and observables}
The orbitals are represented on a uniform grid ($\Delta z=0.3125$, $N=3072$ points in the reference
geometry) with periodic boundary conditions and vacuum margins that prevent overlap across the
boundary. The initial states are obtained by exact diagonalization of the static Hamiltonian with
the Fourier representation of the kinetic energy. They are propagated with the second-order
split-operator method, $\Delta t=0.1$. Since the potential step does not change the density, the
Hartree potential evaluated after each kinetic step is time centered. The Poisson equation is
integrated with a staggered discretization (field at half-integer grid points), which makes the
discrete Coulomb kernel symmetric; a non-symmetric discretization renders the plasma oscillation
of the metal weakly unstable.

The internal field $F(t)$ is the difference of $D(t)g+U_{\rm H}$ between the end of the fifth and the
end of the first unit cell of the crystal divided by their distance, i.e., the field averaged over
0.4 to 2.1~nm from the interface. $F_+$ and $F_-$ are its extrema of the two polarities, and
$F_{\rm dc}$ is its mean over the two central cycles of the pulse. $F_0$ is the peak of $F$ for the
insulated junction ($d=1.5$~nm) at the same $D_0$. The emission spectrum is
$S(\Omega)=|\int dt\,w(t)\dot J(t)e^{i\Omega t}|^2$ with $J=\sum_j\langle\psi_j|p|\psi_j\rangle$ and a
Hann window; harmonic yields $S_n$ are integrals of $S$ over $(n\pm\tfrac12)\omega$. The even-to-odd ratio is $(S_4+S_6)/(S_3+S_5)$. The absorbed energy
is $W=2\sum_k(n_k-n_k^0)E_k/a^2$, where $n_k$ are the occupations of the static eigenstates after
the pulse. The sheet densities quoted in the main text follow from Gauss's law with the effective
permittivity $\varepsilon=D_0/F_0$.

\section{Field dependence}
Table~\ref{tab:field} lists the data underlying Fig.~3(a) and 3(b) of the main text.
\begin{table}[h]
\caption{Field scan for $\Phi_{\rm B}=0.7$~eV, $\ell=8$~nm, wavelength 2.1~$\mu$m, five-cycle pulse.}
\label{tab:field}
\begin{tabular}{ccccccccc}
\hline\hline
$D_0$ & $F_0$ & $F_-$ & $F_+$ & $F_{\rm dc}$ & $W_{\rm contact}$ & $W_{\rm insulated}$ & \multicolumn{2}{c}{even/odd} \\
\multicolumn{5}{c}{(V/nm)} & \multicolumn{2}{c}{(eV/nm$^2$)} & contact & insulated \\
\hline
2.3 & 0.42 & 0.26 & 0.40 & 0.05 & 0.0450 & 0.0006 & 0.07 & 0.018 \\
3.4 & 0.62 & 0.38 & 0.68 & 0.12 & 0.1162 & 0.0007 & 0.23 & 0.032 \\
4.6 & 0.83 & 0.46 & 1.01 & 0.22 & 0.3422 & 0.0019 & 1.48 & 0.038 \\
5.7 & 1.01 & 0.52 & 1.28 & 0.29 & 0.6358 & 0.0108 & 0.86 & 0.039 \\
6.8 & 1.18 & 0.63 & 1.48 & 0.35 & 1.1467 & 0.0504 & 1.35 & 0.037 \\
8.5 & 1.43 & 0.90 & 1.77 & 0.40 & 2.8634 & 0.3088 & 0.92 & 0.024 \\
11.4 & 1.85 & 1.39 & 2.12 & 0.48 & 12.8265 & 2.5270 & 0.94 & 0.098 \\
\hline\hline
\end{tabular}

\end{table}

\section{Convergence}
Table~\ref{tab:conv} shows the dependence of the main observables on the numerical parameters and on
the size of the system for the contacted junction at $D_0=5.7$~V/nm ($F_0=1.0$~V/nm).
\begin{table}[h]
\caption{Convergence tests for the contacted junction at $F_0=1.0$~V/nm. $S_n$ is the yield of the
$n$th harmonic.}
\label{tab:conv}
\begin{tabular}{lcccccccc}
\hline\hline
Run & $F_-$ & $F_+$ & $F_{\rm dc}$ & $W$ & even/odd & $S_4/S_1$ & $S_5/S_1$ & $S_9/S_1$ \\
 & \multicolumn{3}{c}{(V/nm)} & (eV/nm$^2$) & & & & \\
\hline
Reference & 0.52 & 1.28 & $+0.29$ & 0.636 & 0.86 & $2.0\times10^{-4}$ & $2.8\times10^{-4}$ & $1.3\times10^{-5}$ \\
$\Delta t=0.05$ & 0.52 & 1.28 & $+0.29$ & 0.635 & 0.86 & $2.0\times10^{-4}$ & $2.8\times10^{-4}$ & $1.3\times10^{-5}$ \\
$\Delta z=0.234$ & 0.52 & 1.28 & $+0.29$ & 0.640 & 0.87 & $2.0\times10^{-4}$ & $2.8\times10^{-4}$ & $1.3\times10^{-5}$ \\
Crystal 120 cells & 0.52 & 1.28 & $+0.29$ & 0.636 & 0.86 & $2.0\times10^{-4}$ & $2.7\times10^{-4}$ & $1.2\times10^{-5}$ \\
Metal 21 nm & 0.56 & 1.25 & $+0.26$ & 0.540 & 0.87 & $1.9\times10^{-4}$ & $2.5\times10^{-4}$ & $2.3\times10^{-5}$ \\
Metal 34 nm & 0.55 & 1.26 & $+0.27$ & 0.583 & 0.85 & $1.9\times10^{-4}$ & $2.4\times10^{-4}$ & $2.0\times10^{-5}$ \\
\hline\hline
\end{tabular}

\end{table}

\section{Near-field decay length, wavelength, and pulse duration}
Table~\ref{tab:extra} collects additional runs at $D_0=5.7$~V/nm. The crystal without any metal
behaves like the insulated junction. For the contacted junction the rectification
($F_-$, $F_+$, $F_{\rm dc}$) depends only weakly on the decay length $\ell$ of the near field and on
the driving wavelength, and the even-to-odd ratio stays of order unity (0.55 to 1.6), compared with
0.03 to 0.09 without the contact. The quasistatic field and the even-to-odd ratio increase with the
number of cycles, as expected for a charge that accumulates from cycle to cycle. Note that $F$ is
averaged over a fixed window, so that for $\ell=4$~nm the insulated peak value is lower at the same
$D_0$.
\begin{table}[h]
\caption{Additional runs at $D_0=5.7$~V/nm.}
\label{tab:extra}
\begin{tabular}{lcccccccc}
\hline\hline
Run & $F_-$ & $F_+$ & $F_{\rm dc}$ & $W$ & even/odd & $S_4/S_1$ & $S_5/S_1$ & $S_9/S_1$ \\
 & \multicolumn{3}{c}{(V/nm)} & (eV/nm$^2$) & & & & \\
\hline
Contact, $\ell=8$ nm & 0.52 & 1.28 & $+0.29$ & 0.636 & 0.86 & $2.0\times10^{-4}$ & $2.8\times10^{-4}$ & $1.3\times10^{-5}$ \\
Insulated, $\ell=8$ nm & 0.94 & 1.01 & $-0.02$ & 0.011 & 0.04 & $1.8\times10^{-6}$ & $8.1\times10^{-6}$ & $4.8\times10^{-4}$ \\
No metal & 0.94 & 1.02 & $-0.02$ & 0.010 & 0.04 & $1.6\times10^{-6}$ & $7.8\times10^{-6}$ & $4.3\times10^{-4}$ \\
\hline
Contact, $\ell=4$ nm & 0.46 & 1.09 & $+0.25$ & 0.405 & 1.38 & $1.0\times10^{-3}$ & $4.8\times10^{-4}$ & $1.6\times10^{-5}$ \\
Insulated, $\ell=4$ nm & 0.81 & 0.87 & $-0.02$ & 0.001 & 0.03 & $1.5\times10^{-6}$ & $5.4\times10^{-6}$ & $3.1\times10^{-5}$ \\
Contact, $\ell=16$ nm & 0.57 & 1.36 & $+0.31$ & 0.941 & 0.55 & $7.1\times10^{-5}$ & $2.5\times10^{-4}$ & $4.2\times10^{-4}$ \\
Insulated, $\ell=16$ nm & 1.01 & 1.08 & $-0.02$ & 0.046 & 0.09 & $1.9\times10^{-6}$ & $2.2\times10^{-5}$ & $2.2\times10^{-3}$ \\
\hline
Contact, 3.2 $\mu$m & 0.51 & 1.24 & $+0.30$ & 0.825 & 1.56 & $1.8\times10^{-4}$ & $5.6\times10^{-5}$ & $8.8\times10^{-4}$ \\
Insulated, 3.2 $\mu$m & 0.95 & 1.02 & $-0.02$ & 0.007 & 0.03 & $1.7\times10^{-6}$ & $2.3\times10^{-6}$ & $1.5\times10^{-6}$ \\
\hline
Contact, 8 cycles & 0.53 & 1.32 & $+0.38$ & 1.172 & 1.47 & $2.5\times10^{-4}$ & $2.5\times10^{-4}$ & $1.2\times10^{-4}$ \\
Insulated, 8 cycles & 0.99 & 1.01 & $-0.01$ & 0.016 & 0.01 & $1.9\times10^{-7}$ & $7.3\times10^{-6}$ & $6.4\times10^{-4}$ \\
\hline\hline
\end{tabular}

\end{table}

\section{Far-field parity of the contact emission}
Consider an antenna with mirror symmetry $x\to-x$ driven at its dipolar resonance. Let $\hat n_{\rm R}$
and $\hat n_{\rm L}=-\hat n_{\rm R}$ be the outward in-plane normals of the two ends. The in-plane near
field has the same sign at both ends, so the field component along the outward normal obeys
$E_{\rm L}(t)=-E_{\rm R}(t)=E_{\rm R}(t+T/2)$ for a periodic driver: the left end is driven like the
right end half a period later. The current along the outward normal at the left end is therefore
$j(t+T/2)$ if $j(t)$ is that at the right end, and the in-plane currents are
$J_{x,{\rm R}}=j(t)$ and $J_{x,{\rm L}}=-j(t+T/2)$. The $n$th Fourier component of their sum is
proportional to $1-e^{-in\pi}=1-(-1)^n$, so even harmonics cancel in the direction normal to the
surface and odd harmonics add. This holds for any source that depends only on the local field
relative to the outward normal, including the field-gradient mechanism. Even harmonics survive off
axis, in odd diffraction orders of an array, and for antennas without mirror symmetry.
\end{document}
